\clearpage
\chap{The residual floor and the closure of admissible changes}\label{sec:floor}
An algebraic obstruction asks whether a residual is exactly the defect of an allowed correction. A numerical floor asks how closely allowed defects can approximate it. These questions agree under a closedness hypothesis. Without that hypothesis they can disagree even in a complete Hilbert setting.

\subsection{Distance to the actual correction image}
Let \(R\) be a normed additive group and \(B\le R\) an additive subgroup. In the protected linear model, \(B=d(\ker p)\). For a residual \(\Omega\in R\), define
\[
 \nu_B(\Omega)=\inf_{b\in B}\norm{\Omega-b}.
\]
The subgroup contains zero, so this is a finite nonnegative real number. It is an infimum, not an assertion that a nearest admissible residual exists.

\begin{proposition}[The floor detects closure]\label{KS76}
For every \(b\in B\),
\[
 0\le\nu_B(\Omega)\le\norm{\Omega-b}.
\]
Moreover,
\[
 \nu_B(\Omega)=0\quad\Longleftrightarrow\quad\Omega\in\overline B.
\]
If \(B\) is closed, this is equivalent to \(\Omega\in B\).
\end{proposition}
\begin{proof}
The set of candidate distances is nonempty because \(0\in B\), bounded below by zero, and contains the finite value \(\norm\Omega\). Its infimum is therefore a finite nonnegative real, and the defining lower-bound property gives \(\nu_B(\Omega)\le\norm{\Omega-b}\) for every \(b\in B\).

Suppose \(\nu_B(\Omega)=0\) and fix \(\varepsilon>0\). There must be \(b\in B\) with \(\norm{\Omega-b}<\varepsilon\). Otherwise every candidate distance would be at least \(\varepsilon\), making \(\varepsilon\) a lower bound and forcing the infimum to be at least \(\varepsilon\), a contradiction. Thus every open ball about \(\Omega\) meets \(B\), which is the closure criterion.

Conversely, if \(\Omega\in\overline B\), every such ball supplies \(b\) with \(\norm{\Omega-b}<\varepsilon\). The infimum then lies between zero and \(\varepsilon\) for every positive \(\varepsilon\), so it is zero. If \(B\) is closed, membership in its closure is exactly membership in \(B\). No step asserts that the infimum has a minimizing element when \(B\) is not closed, and closedness alone is being used here only for the zero-floor membership conclusion.
\end{proof}
\source{KS76}
Translation by a member \(b_0\in B\) does not change the floor: replacing \(b\) by \(b-b_0\) gives the same set of distances. Thus the floor is well defined on the algebraic quotient \(R/B\). When \(B\) is a linear subspace, this is the usual quotient seminorm. A nonzero algebraic class can have seminorm zero when \(B\) is not closed. The retained source implements the floor on residual representatives and proves the closure criterion. This quotient interpretation is its elementary translation-invariance calculation, not a new claim that every obstruction quotient carries a separated norm.

\subsection{A precise three-way distinction in this fixed model}
For this specified subgroup and residual there are exactly three possibilities:
\begin{center}
\begin{tabularx}{\linewidth}{@{}p{.29\linewidth}p{.20\linewidth}X@{}}
\toprule Residual location & Algebraic class & Meaning of the floor\\\midrule
\(\Omega\in B\) & Zero & Exact cancellation exists, and the floor is zero.\\
\(\Omega\in\overline B\setminus B\) & Nonzero & Arbitrarily small corrected residuals exist, but no exact cancellation exists.\\
\(\Omega\notin\overline B\) & Nonzero & The floor is strictly positive and blocks every allowed correction.\\\bottomrule
\end{tabularx}
\end{center}
This is a classification for the stated normed-group correction problem, not a universal trichotomy for every possible Komorebi model. Testing which case holds can itself be difficult. The definitions do not supply a decision algorithm.

In finite-dimensional real or complex linear models, every linear subspace is closed. The middle row is therefore absent. It becomes possible in infinite dimensions, or for additive subgroups that are not closed real-linear subspaces. For example the rational subgroup of \(\R\) is dense, so the distance from any irrational real to that subgroup is zero even though its coset is nonzero.

Closedness also should not be confused with a numerical error bound back in the candidate space. The statement \(\Omega\in B\) supplies an admissible correction, but not a bound on its size. The closed-range estimates of Sections~\ref{sec:closed} and~\ref{sec:motion}, or their protected counterparts, are the additional analytic tools that control that size in their declared settings.

\subsection{A persistent floor blocks every correction sequence}
The correction image can vary with the resolution or the admissibility rule. Let \(B_n\le R\) and let \(\Omega_n\in R\) be given.

\begin{proposition}[Uniformly positive floors obstruct asymptotic repair]\label{KS77}
Suppose \(c>0\) and eventually
\[
 c\le\nu_{B_n}(\Omega_n).
\]
For any choices \(b_n\in B_n\), it is impossible that
\[
 \norm{\Omega_n-b_n}\longrightarrow0.
\]
\end{proposition}
\begin{proof}
Every allowed correction satisfies
\(\norm{\Omega_n-b_n}\ge\nu_{B_n}(\Omega_n)\). The corrected norms are therefore eventually at least \(c\). They cannot converge to zero.
\end{proof}
\source{KS77}
The quantifier covers all admissible sequences, not just the output of one repair routine. On the other hand, the hypothesis is an eventually \emph{uniform} positive lower bound. Having a nonzero obstruction class at each stage is not the same statement. The classes might have zero floors, or their positive floors might tend to zero.

A finite-dimensional illustration uses \(R=\R^2\), \(B=\R\times\{0\}\), and \(\Omega_n=(n,1)\). Every correction leaves the second coordinate equal to one, so the floor is one independently of \(n\). Replacing the residuals by \((n,1/(n+1))\) gives nonzero algebraic classes at every stage, but their floors tend to zero. The persistent-floor theorem applies to the first sequence and not the second.

\subsection{The bridge back to exactification}
Theorem~\ref{KS73} and Proposition~\ref{KS76} fit together closely. If corrections produce a norm-convergent, protected path with vanishing residual, the actual quotient class vanishes. If only the corrected residuals tend to zero, all that follows is membership in \(\overline{d(\ker p)}\).

In particular, when the class is nonzero but its floor is zero, \emph{no} protected path with vanishing residual can satisfy all the exactification hypotheses. Something must fail: the path may escape, its movements may not be summable, the candidate space may not be complete, or its residual and semantic maps may not have the required continuity. The next model isolates escape while keeping the spaces Hilbert and both maps continuous.
